\chapter{Лекция}
\label{ch:intro}

\section{Архитектура приемника}

Рассмотрим архитектуру простого приемника

\begin{figure}[H]
    \centering
    \includegraphics[width=1.0\textwidth]{rx.png}
    \caption{Архитектура простого приемника}
\end{figure}

\subsection{Радиоканал и SDR}

Из радиоканала сигнал попадает на антенну, где преобразуется в электрический сигнал и попадает в SDR, которая дискретизирует сигнал, производит
какие-то дополнительные манипуляции и на выходе выдает отсчеты сигнала. 

\subsection{Matched filter}

Далее семплы попадают на согласованный фильтр (matched filter), который позволяет увеличить SNR (Signal Noise Ration), т.е выделить
полезный сигнал из общего сигнала с шумом. Импульсная характеристика фильтра должна согласовываться с импульсной характеристикой 
фильтра в передатчике. Допустим, передатчик имеет импульсную характеристику $g(n)$, тогда на приемной стороне ИХ фильтра будет следующей:

$$h(n) = Kg*(L-n)$$

ИХ согласованного фильтра на приемнике - ИХ, которая является комплексно-сопряженной (зеркальной) с ИХ передатчика и сдвинутой на L отсчетов, где
L - кол-во семплов на символ. $K$ - коэффициент.\\

Мы используем прямоугольную ИХ на передатчике, поэтому на приеме эта ИХ никак не изменится. \\

На выходе фильтра получаем пилообразный график:

\begin{figure}[H]
    \centering
    \includegraphics[width=1.0\textwidth]{amatch.png}
    \caption{Выход matched filter}
\end{figure}

\subsection{Downsampling}
На стороне передатчика мы делали операцию upsampling, которая повышала частоту дискретизации. Это позволяло нам сделать так, чтобы на
1 символ приходилось 10 семплов. На стороне приема делается обратная операция - downsampling, т.е нам нужно из каждых 10 семплов
выбрать один такой, который наиболее точно описывает переданный символ. Но какой семпл нужно выбирать? После прохождения через фильтр
каждый 10-ый семпл будет иметь наибольший SNR, именно его и нужно выбирать. Но с какого семпла нам нужно начинать выполнять эту операцию? Ведь
мы можем "промахнуться" и начать выбирать не те семплы. На этот вопрос отвечает символьная синхронизация, но это тема последующих
занятий. На этом занятии будет начинать downsampling с первого семпла.

\subsection{Quantization}

Проходя через радиоканал, сигнал искажается и принятые символы будут искаженными. Если на передаче отправлялись символы \{1,0,1,0,-1,0\},
то на приеме можем получить \{0.74,-0.32,1.24,0.2,-1.35,0.01\}. Четко определить, какой последовательности битов соотвествуют символы невозможно,
поэтому перед тем, как начать этап преобразования символов в биты, нужно понять, к какой из возможных точек созвездия эти точки находятся
ближе всего, и как-бы округлить координаты точек так, чтобы они точно соотвествовали координатам истинных точек созвездия. Чтобы оценивать, к какой
истинной точке на созвездии ближе всего находится точка с искаженными координатами, я буду вычислять расстояние между искаженной точкой и всеми
истинными точками. \\

Пример вычисления расстояния между точками сигнального созвездия BPSK:

\begin{figure}[H]
    \centering
    \includegraphics[width=1.0\textwidth]{dist.png}
    \caption{Расстояние между точками созвездия}
\end{figure}

В созвездии BPSK имеются 2 точки с координатами (-1,0) и (1,0), точка A имеет искаженные координаты. По теореме Пифагора можно найти
расстояния между точками. Точка A ближе находится к точке (1,0), поэтому логично предположить, что передавалась единица, а не ноль.

\subsection{Demapper}

На передаче мы подавали биты на специальное устройство под названием Mapper, которое ставило в соответствие последовательности бит
символы. Demapper делает обратную задачу - символам сопоставляет биты.

\subsection{Decoder}

На стороне передатчика передается текст, поэтому на стороне приема нужен блок, который преобразовывал бы биты в текст. Этой задачей
будет заниматься блок Decoder.

\endinput
